% Partie : counterfactuals
% Chapitre : minimal-change-semantics
% Section : introduction

\documentclass[../../../include/open-logic-section]{subfiles}

\begin{document}

\olfileid[fr]{cnt}{min}{int}

\olsection{Introduction}

Stalnaker et Lewis ont proposé des analyses de conditionnels
contrefactuels tels que « Si l'on frottait cette allumette, elle
s'enflammerait. » Leurs propositions visaient à préciser les
conditions de vérité de tels énoncés. L'idée qui leur est commune est
la suivante : pour évaluer un conditionnel contrefactuel, il faut
considérer les mondes possibles qui diffèrent le moins du monde actuel
tout en rendant l'antécédent vrai. Si le conséquent est vrai dans ces
mondes, le contrefactuel est vrai. Supposons, par exemple, que je
tienne une allumette et un carnet d'allumettes dans ma main. Dans le
monde actuel, je me contente de les regarder et de me demander ce
qui se passerait si je frottais l'allumette. Le changement minimal qui
conduit à un monde où je la frotte consiste à décider de le faire,
puis à la frotter. Il est minimal parce que rien d'autre ne change :
je ne saute pas en l'air, le frottement de l'allumette n'enflamme pas
aussi mes cheveux, mes doigts ne perdent pas soudain toute leur force
et personne ne m'asperge d'eau par surprise avec un pistolet à eau
Super Soaker. Dans cette autre situation possible, l'allumette
s'enflamme. Il est donc vrai que, si je la frottais, elle
s'enflammerait.

Cette idée intuitive peut recevoir une sémantique formelle pour les
logiques des conditionnels contrefactuels. Lewis a introduit le
symbole « $\boxright$ » pour le contrefactuel, tandis que Stalnaker
employait « $>$ ». Nous utiliserons~$\cif$ et l'ajouterons comme
connecteur binaire à la logique propositionnelle. Ainsi, en plus des
!!{formula}s de la forme $!A \lif !B$, nous aurons des !!{formula}s
de la forme~$!A \cif !B$. La sémantique formelle, comme la sémantique
relationnelle de la logique modale, repose sur des modèles dans
lesquels les !!{formula}s sont évaluées en des mondes. La condition
de satisfaction qui définit $\mSat{M}{!A \cif !B}[w]$ fait
intervenir $\mSat{M}{!A}[w']$ et $\mSat{M}{!B}[w']$ pour certains
autres mondes~$w'$. De quels mondes~$w'$ s'agit-il ? Intuitivement,
ceux qui sont les plus proches de~$w$ parmi les mondes où
$\mSat{M}{!A}[w']$ vaut. Le modèle
doit donc également comporter une relation de « proximité ».

Lewis a proposé une représentation graphique instructive des
situations contrefactuelles. Chaque monde possible se trouve au
centre d'un ensemble de sphères emboîtées contenant d'autres mondes :
nous les dessinons sous forme de cercles concentriques. Les mondes
situés entre deux sphères sont aussi proches les uns que les autres
du monde central ; ceux d'une sphère intérieure en sont plus proches,
et ceux d'une sphère extérieure en sont plus éloignés.
\begin{center}
\begin{tikzpicture}[scale=.7]\small
  \spheresystem{5}
  \draw (0,0) node {$w$};
  \propositionintersect{0}{4}{40}{3.5}
  {\tikzset{proposition/.append={smooth,tension=1.4}}
  \proposition[shift={(1.6,-1)}]{90}{1}{30}{4}}
  \path (1.5,3) node[below] {$\formula{B}$};
  \path (3,0) node {$\formula{A}$};
\end{tikzpicture}
\end{center}
Les mondes-$!A$ les plus proches sont les mondes~$w'$ où~$!A$ est vraie
et qui se trouvent dans la plus petite sphère autour du monde
central~$w$ qui contient de tels mondes (zone grise). Intuitivement,
$!A \cif !B$ est vrai en~$w$ si~$!B$ est vrai dans tous les mondes-$!A$
les plus proches.

\end{document}
